\title{Geometric Interpretation of Voltage Constraints in Permanent-Magnet Synchronous Machines}
\author{}
\date{24 September 2026}
\newcommand{\PaperAbstract}{%
This paper derives the steady-state voltage limit of a permanent-magnet
synchronous machine as an affine quadratic inequality in the $dq$ current
plane.  The conic type is determined from the quadratic part rather than
assumed from machine terminology.  Completing the square identifies the
center; spectral decomposition identifies the principal directions and
semi-axis lengths; and the classical double-angle rotation formula is shown
to be the same calculation in coordinates.  Parameter limits clarify the
roles of resistance, speed, saliency, permanent-magnet flux, and available
voltage.  Current and torque constraints are then added, so maximum torque per
ampere, field weakening, and maximum torque per voltage emerge as tangency
problems.  Degenerate and non-elliptic cases are stated explicitly.}
